\documentclass[10pt,a4paper]{amsart}
\usepackage{amsmath,amssymb,amsthm,graphicx}
\usepackage{verbatim}
\newenvironment{reference}{\comment}{\endcomment}
%\newenvironment{reference}{}{}
\newenvironment{slogan}{\comment}{\endcomment}
\newenvironment{history}{\comment}{\endcomment}

% For commutative diagrams we use Xy-pic
\usepackage[all]{xy}

% We use 2cell for 2-commutative diagrams.
\xyoption{2cell}
% Theorem environments.
%
\theoremstyle{plain}
\newtheorem{theorem}[subsection]{Theorem}
\newtheorem{proposition}[subsection]{Proposition}
\newtheorem{lemma}[subsection]{Lemma}

\theoremstyle{definition}
\newtheorem{definition}[subsection]{Definition}
\newtheorem{example}[subsection]{Example}
\newtheorem{exercise}[subsection]{Exercise}
\newtheorem{situation}[subsection]{Situation}

\theoremstyle{remark}
\newtheorem{remark}[subsection]{Remark}
\newtheorem{remarks}[subsection]{Remarks}

\numberwithin{equation}{subsection}

% Macros
%
\def\lim{\mathop{\mathrm{lim}}\nolimits}
\def\colim{\mathop{\mathrm{colim}}\nolimits}
\def\Spec{\mathop{\mathrm{Spec}}}
\def\Hom{\mathop{\mathrm{Hom}}\nolimits}
\def\Ext{\mathop{\mathrm{Ext}}\nolimits}
\def\SheafHom{\mathop{\mathcal{H}\!\mathit{om}}\nolimits}
\def\SheafExt{\mathop{\mathcal{E}\!\mathit{xt}}\nolimits}
\def\Sch{\mathit{Sch}}
\def\Mor{\mathop{\mathrm{Mor}}\nolimits}
\def\Ob{\mathop{\mathrm{Ob}}\nolimits}
\def\Sh{\mathop{\mathit{Sh}}\nolimits}
\def\NL{\mathop{N\!L}\nolimits}
\def\CH{\mathop{\mathrm{CH}}\nolimits}
\def\proetale{{pro\text{-}\acute{e}tale}}
\def\etale{{\acute{e}tale}}
\def\QCoh{\mathit{QCoh}}
\def\Ker{\mathop{\mathrm{Ker}}}
\def\Im{\mathop{\mathrm{Im}}}
\def\Coker{\mathop{\mathrm{Coker}}}
\def\Coim{\mathop{\mathrm{Coim}}}

% Boxtimes
%
\DeclareMathSymbol{\boxtimes}{\mathbin}{AMSa}{"02}

%
% Macros for moduli stacks/spaces
%
\def\QCohstack{\mathcal{QC}\!\mathit{oh}}
\def\Cohstack{\mathcal{C}\!\mathit{oh}}
\def\Spacesstack{\mathcal{S}\!\mathit{paces}}
\def\Quotfunctor{\mathrm{Quot}}
\def\Hilbfunctor{\mathrm{Hilb}}
\def\Curvesstack{\mathcal{C}\!\mathit{urves}}
\def\Polarizedstack{\mathcal{P}\!\mathit{olarized}}
\def\Complexesstack{\mathcal{C}\!\mathit{omplexes}}
% \Pic is the operator that assigns to X its picard group, usage \Pic(X)
% \Picardstack_{X/B} denotes the Picard stack of X over B
% \Picardfunctor_{X/B} denotes the Picard functor of X over B
\def\Pic{\mathop{\mathrm{Pic}}\nolimits}
\def\Picardstack{\mathcal{P}\!\mathit{ic}}
\def\Picardfunctor{\mathrm{Pic}}
\def\Deformationcategory{\mathcal{D}\!\mathit{ef}}
\usepackage{cleveref}
\crefname{originallemma}{original lemma}{original lemmas}
\Crefname{originallemma}{Original lemma}{Original lemmas}
\makeatletter
\newlabel{lemma-cocontinuous-morphism-topoi}{{00XO}{source}{}{sourcepointer-1}{}}
\newlabel{lemma-cocontinuous-morphism-topoi@cref}{{[originallemma][0][]00XO}{[1][0][]source}{}{}{}}
\makeatother
\crefname{originallemma}{original lemma}{original lemmas}
\Crefname{originallemma}{Original lemma}{Original lemmas}
\makeatletter
\newlabel{lemma-when-shriek}{{00XR}{source}{}{sourcepointer-2}{}}
\newlabel{lemma-when-shriek@cref}{{[originallemma][0][]00XR}{[1][0][]source}{}{}{}}
\makeatother
\title{A cocontinuous and continuous functor satisfying five conditions induces an equivalence of topoi}
\author{The Stacks Project Authors}
\date{stacks-project:sites.tex:local-snapshot-20261012; GNU FDL 1.2 or later}
\begin{document}
\maketitle
\noindent\textit{Editorial note.} All delivered mathematical text is the author's own, copied from the pinned Stacks Project chapter file sites-cohomology.tex; no statement or proof step was rewritten, reordered or removed.\par\medskip
\noindent\textit{Editorial note.} The proof is the author's complete argument; every step the author writes out is inside the counted span and no closing sentence is a placeholder.\par\medskip
\noindent\textit{Editorial note.} Every labelled reference occurring inside the counted spans is carried in assembly.external\_references as a non-load-bearing navigation pointer (same-chapter ones at their definition site, cross-chapter ones at the invocation inside the counted span), so no load-bearing mathematics is replaced by a stub. No bibliography citation occurs inside the counted spans.\par\medskip
\noindent\textit{Editorial note.} The checked chapter list for same-name collisions is the set of stacks chapters staged in this repository: chow, cohomology, divisors, etale-cohomology, flat, more-algebra, more-morphisms, morphisms, spaces-more-morphisms, varieties together with this card's chapter sites-cohomology.\par\medskip
\noindent\textit{Editorial note.} The preamble is the chapter's own retained minimal preamble (sources/stacks-sites-cohomology/retained-preamble.tex), a verbatim excerpt of the native preamble: the Xy-pic setup and the native theorem declarations with the mathematical macros.\par\medskip
\noindent\textit{Allowed standard prerequisites (editorial).}\begin{itemize}
\item The topos morphism associated to a cocontinuous functor. Same chapter, Lemma lemma-cocontinuous-morphism-topoi. Supplies the morphism of topoi whose adjunction maps are shown to be isomorphisms.
\item The pullback along a cocontinuous functor and its value on objects. Same chapter, Lemma lemma-when-shriek. Identifies the pullback of a sheaf with the value at the image object.
\item The pushforward as a limit over morphisms into an object. Same chapter, the description of the pushforward of a cocontinuous functor as a limit over the category of morphisms from the image of an object. Supplies the description of the composite of pullback and pushforward used in the proof.
\end{itemize}\medskip
\ifdefined\hypertarget\hypertarget{sourcepointer-1}{}\fi
\noindent\textit{Source reference.} Original lemma 00XO: sources/stacks-sites/pinned-original-source.tex, source lines 7013--7013. Navigation pointer invoked at proof line 7013; the counted span contains the pointer only and the cited result is not reproduced inside it, so the pointer is registered as a non-load-bearing navigation stub in the established form for Stacks material. No bibliography citation is involved.\par\medskip
\ifdefined\hypertarget\hypertarget{sourcepointer-2}{}\fi
\noindent\textit{Source reference.} Original lemma 00XR: sources/stacks-sites/pinned-original-source.tex, source lines 6954--6954. Navigation pointer invoked at proof line 6954; the counted span contains the pointer only and the cited result is not reproduced inside it, so the pointer is registered as a non-load-bearing navigation stub in the established form for Stacks material. No bibliography citation is involved.\par\medskip
\begin{lemma}
\label{lemma-equivalence}
Let $\mathcal{C}$, $\mathcal{D}$ be sites.
Let $u : \mathcal{C} \to \mathcal{D}$ be a functor.
Assume that
\begin{enumerate}
\item $u$ is cocontinuous,
\item $u$ is continuous,
\item given $a, b : U' \to U$ in $\mathcal{C}$ such that
$u(a) = u(b)$, then there exists a covering $\{f_i : U'_i \to U'\}$
in $\mathcal{C}$ such that $a \circ f_i = b \circ f_i$,
\item given $U', U \in \Ob(\mathcal{C})$ and
a morphism $c : u(U') \to u(U)$ in $\mathcal{D}$ there exists
a covering $\{f_i : U_i' \to U'\}$ in $\mathcal{C}$
and morphisms $c_i : U_i' \to U$ such that $u(c_i) = c \circ u(f_i)$, and
\item given $V \in \Ob(\mathcal{D})$ there exists a covering
of $V$ in $\mathcal{D}$ of the form $\{u(U_i) \to V\}_{i \in I}$.
\end{enumerate}
Then the morphism of topoi
$$
g : \Sh(\mathcal{C}) \longrightarrow \Sh(\mathcal{D})
$$
associated to the cocontinuous functor $u$ by
Lemma \ref{lemma-cocontinuous-morphism-topoi}
is an equivalence.
\end{lemma}

\begin{proof}
Assume $u$ satisfies properties (1) -- (5). We will show that
the adjunction mappings
$$
\mathcal{G} \longrightarrow g_*g^{-1}\mathcal{G}
\quad\text{and}\quad
g^{-1}g_*\mathcal{F} \longrightarrow \mathcal{F}
$$
are isomorphisms.

\medskip\noindent
Note that Lemma \ref{lemma-when-shriek} applies and we have
$g^{-1}\mathcal{G}(U) = \mathcal{G}(u(U))$ for any sheaf $\mathcal{G}$
on $\mathcal{D}$. Next, let $\mathcal{F}$ be a sheaf on $\mathcal{C}$,
and let $V$ be an object of $\mathcal{D}$. By definition we have
$g_*\mathcal{F}(V) = \lim_{u(U) \to V} \mathcal{F}(U)$.
Hence
$$
g^{-1}g_*\mathcal{F}(U) = \lim_{U', u(U') \to u(U)} \mathcal{F}(U')
$$
where the morphisms $\psi : u(U') \to u(U)$ need not be of the form
$u(\alpha)$. The category of such pairs $(U', \psi)$ has a final
object, namely $(U, \text{id})$, which gives rise to the map from
the limit into $\mathcal{F}(U)$. Let $(s_{(U', \psi)})$ be an element
of the limit. We want to show that $s_{(U', \psi)}$ is uniquely determined
by the value $s_{(U, \text{id})} \in \mathcal{F}(U)$. By property (4) given
any $(U', \psi)$ there exists a covering $\{U'_i \to U'\}$ such that the
compositions $u(U'_i) \to u(U') \to u(U)$ are of the form $u(c_i)$
for some $c_i : U'_i \to U$ in $\mathcal{C}$. Hence
$$
s_{(U', \psi)}|_{U'_i} = c_i^*(s_{(U, \text{id})}).
$$
Since $\mathcal{F}$ is a sheaf it follows that indeed $s_{(U', \psi)}$
is determined by $s_{(U, \text{id})}$. This proves uniqueness.
For existence, assume given any
$s \in \mathcal{F}(U)$, $\psi : u(U') \to u(U)$, $\{f_i : U_i' \to U'\}$
and $c_i : U_i' \to U$ such that $\psi \circ u(f_i) = u(c_i)$ as above.
We claim there exists a (unique) element
$s_{(U', \psi)} \in \mathcal{F}(U')$ such that
$$
s_{(U', \psi)}|_{U'_i} = c_i^*(s).
$$
Namely, a priori it is not clear the elements
$c_i^*(s)|_{U_i' \times_{U'} U_j'}$
and $c_j^*(s)|_{U_i' \times_{U'} U_j'}$ agree, since
the diagram
$$
\xymatrix{
U_i' \times_{U'} U_j' \ar[r]_-{\text{pr}_2} \ar[d]_{\text{pr}_1} &
U_j' \ar[d]^{c_j} \\
U_i' \ar[r]^{c_i} & U}
$$
need not commute. But condition (3) of the lemma guarantees that there
exist coverings
$\{f_{ijk} : U'_{ijk} \to U_i' \times_{U'} U_j'\}_{k \in K_{ij}}$ such that
$c_i \circ \text{pr}_1 \circ f_{ijk} = c_j \circ \text{pr}_2 \circ f_{ijk}$.
Hence
$$
f_{ijk}^* \left(c_i^*s|_{U_i' \times_{U'} U_j'}\right)
=
f_{ijk}^* \left(c_j^*s|_{U_i' \times_{U'} U_j'}\right)
$$
Hence $c_i^*(s)|_{U_i' \times_{U'} U_j'} = c_j^*(s)|_{U_i' \times_{U'} U_j'}$
by the sheaf condition for $\mathcal{F}$ and hence the existence of
$s_{(U', \psi)}$ also by the sheaf condition for $\mathcal{F}$. The uniqueness
guarantees that the collection $(s_{(U', \psi)})$ so obtained is an element
of the limit with $s_{(U, \psi)} = s$. This proves
that $g^{-1}g_*\mathcal{F} \to \mathcal{F}$ is an isomorphism.

\medskip\noindent
Let $\mathcal{G}$ be a sheaf on $\mathcal{D}$. Let $V$ be an
object of $\mathcal{D}$. Then we see that
$$
g_*g^{-1}\mathcal{G}(V) =
\lim_{U, \psi : u(U) \to V} \mathcal{G}(u(U))
$$
By the preceding paragraph we see that the value of the sheaf
$g_*g^{-1}\mathcal{G}$ on an object $V$ of the form $V = u(U)$
is equal to $\mathcal{G}(u(U))$. (Formally, this holds because
we have $g^{-1}g_*g^{-1} \cong g^{-1}$, and the description
of $g^{-1}$ given at the beginning of the proof; informally just by
comparing limits here and above.)
Hence the adjunction mapping $\mathcal{G} \to g_*g^{-1}\mathcal{G}$ has
the property that it is a bijection on sections over any object of the
form $u(U)$. Since by axiom (5) there
exists a covering of $V$ by objects of the form $u(U)$ we see
easily that the adjunction map is an isomorphism.
\end{proof}

\end{document}
